\begin{proposition}[切向量的两种定义]
设 $M$ 为光滑流形、$p\in M$。通过 $p$ 的光滑曲线按相同初速度取商所得空间，与点 $p$ 处满足 Leibniz 规则的实线性导子空间自然同构。对应为 $[\gamma]\mapsto[f\mapsto (f\circ\gamma)^\prime(0)]$。
\end{proposition}
\begin{proof}
需证两个方向：每个曲线速度是导子，每个导子由某曲线速度实现。

方向一：曲线速度 $\Rightarrow$ 导子。任给光滑曲线 $\gamma:(-\varepsilon,\varepsilon)\to M$，$\gamma(0)=p$，定义 $v_\gamma(f)=\left.\frac{\dd}{\dd t}\right|_0 f(\gamma(t))$。对 $f,g\in C^\infty(M)$，由链式法则
\begin{align*}
 v_\gamma(fg)
 &=\left.\frac{\dd}{\dd t}\right|_0\bigl(f(\gamma(t))\,g(\gamma(t))\bigr)\\
 &=\dot f(\gamma(0))\,g(\gamma(0))+f(\gamma(0))\,\dot g(\gamma(0))\\
 &=v_\gamma(f)\,g(p)+f(p)\,v_\gamma(g),
\end{align*}
故 $v_\gamma$ 满足 Leibniz 规则，是 $T_pM$ 中的导子。线性性 $v_\gamma(af+bg)=a\,v_\gamma(f)+b\,v_\gamma(g)$ 直接由导数线性得到。

方向二：导子 $\Rightarrow$ 曲线速度。先说明导子的坐标表示没有额外假设。Leibniz 规则给 $v(1)=0$。若函数 $f$ 在 $p$ 附近恒零，取邻域内等于一且支撑在该邻域的光滑截断函数 $\chi$，则 $\chi f=0$ 且
\[
v(\chi f)=v(\chi)f(p)+\chi(p)v(f)=v(f)=0.
\]
因此 $v$ 只依赖函数在 $p$ 的芽。坐标邻域内的 Hadamard 分解给
\[
f(x)-f(p)=\sum_i(x^i-x^i(p))h_i(x),\qquad h_i(p)=\partial_i f(p).
\]
对它作用 $v$，由 Leibniz 规则得到 $v(f)=\sum_i v(x^i)\partial_i f(p)$。

给定 $v\in T_pM$（导子定义），在局部坐标 $(x^1,\ldots,x^n)$ 中写
\begin{equation*}
 v=v^i\,\partial_i\big|_p,
 \qquad
 v^i=v(x^i).
\end{equation*}
构造坐标曲线 $\gamma:(-\varepsilon,\varepsilon)\to M$，令
\begin{equation*}
 x^i(\gamma(t))=x^i(p)+t\,v^i,
 \qquad i=1,\ldots,n,
\end{equation*}
在足够小 $\varepsilon$ 上限制于坐标邻域内。对任意 $f\in C^\infty(M)$，在坐标中写 $f=f(x^1,\ldots,x^n)$，则
\begin{align*}
 \left.\frac{\dd}{\dd t}\right|_0 f(\gamma(t))
 &=\left.\frac{\dd}{\dd t}\right|_0 f\bigl(x^1(p)+tv^1,\ldots,x^n(p)+tv^n\bigr)\\
 &=\sum_{i=1}^n \frac{\partial f}{\partial x^i}\bigg|_p v^i
 =v^i\,\partial_i f(p)=v(f).
\end{align*}
因此 $v_\gamma=v$，每个导子都由某条曲线速度实现。

若两条曲线 $\gamma_1,\gamma_2$ 给出同一导子，则对所有 $f$ 有 $v_{\gamma_1}(f)=v_{\gamma_2}(f)$，特别取 $f=x^i$ 得 $\dot x^i_1(0)=\dot x^i_2(0)$，即速度分量相同。反之，速度分量相同则对任意 $f$ 的一阶作用相同。因此映射 $[\gamma]\mapsto v_\gamma$ 是曲线等价类到导子的双射，且与坐标选取无关，两种定义自然同构。
\end{proof}

\begin{lemma}[Poincaré 引理与缩放同伦]
设 $U\subset\RR^n$ 关于原点星形。对 $k\ge1$ 定义 $(K\alpha)_x(v_1,\ldots,v_{k-1})=\int_0^1t^{k-1}\alpha_{tx}(x,v_1,\ldots,v_{k-1})\,\dd t$。则 $\dd K+K\dd=\id-F_0^*$；特别地，每个正次数闭形式 $\alpha$ 都满足 $\alpha=\dd K\alpha$。
\end{lemma}
\begin{proof}
用缩放同伦 $F_t(x)=tx$ 将 $t=1$ 的形式 $\alpha$ 连续拉回到 $t=0$（原点），沿途的"变化率"由 Cartan 公式控制，积分后得到同伦恒等式。

对 $k$-形式 $\alpha$，定义
\begin{equation*}
 H:[0,1]\times U\to U,
 \qquad H(t,x)=tx.
\end{equation*}
沿 $t$ 方向的生成向量场在 $U$ 中对应径向场 $\mathcal R/t$（因 $\frac{\dd}{\dd t}(tx)=x=\mathcal R/t$ 在 $tx$ 处取值）。

由 Lie 导数与外微分的关系（Cartan 公式）$\mathcal L_X=\dd\iota_X+\iota_X\dd$ 以及拉回与 Lie 导数的交换关系，
\begin{align*}
 \frac{\dd}{\dd t}F_t^*\alpha
 &=F_t^*(\mathcal L_{\mathcal R/t}\alpha)
 =F_t^*\left(\dd\iota_{\mathcal R/t}\alpha+\iota_{\mathcal R/t}\dd\alpha\right)\\
 &=\frac{1}{t}F_t^*\left(\dd\iota_{\mathcal R}\alpha+\iota_{\mathcal R}\dd\alpha\right).
\end{align*}

从 $t=0$ 积分到 $1$：
\begin{equation*}
 \alpha-F_0^*\alpha
 =\int_0^1\frac{\dd}{\dd t}F_t^*\alpha\,\dd t
 =\dd\!\int_0^1 F_t^*(\iota_{\mathcal R}\alpha)\,\frac{\dd t}{t}
 +\int_0^1 F_t^*(\iota_{\mathcal R}\dd\alpha)\,\frac{\dd t}{t}.
\end{equation*}
拉回包含切向量的缩放，不能把径向向量在 $tx$ 处的值再重复乘一次 $t$。对 $k$-形式取分量后得到 $t^{k-1}\alpha_{tx}(x,\cdot)$；不要求 $\alpha$ 本身齐次。记
\begin{equation*}
 (K\alpha)_x=\int_0^1 F_t^*(\iota_{\mathcal R}\alpha)\,\frac{\dd t}{t},
\end{equation*}
即\eqnrefs{eq:poincare-homotopy-operator}，则上式化为
\begin{equation*}
 \alpha-F_0^*\alpha=\dd K\alpha+K\dd\alpha,
 \end{equation*}
也就是\eqnrefs{eq:poincare-homotopy-identity}。

若 $k\ge1$ 且 $\dd\alpha=0$（闭形式），常值映射 $F_0$ 把 $k$-形式拉回为零（$k\ge1$ 时原点处切空间退化），故 $F_0^*\alpha=0$。代入同伦恒等式：
\begin{equation*}
 \alpha=\dd(K\alpha),
\end{equation*}
即 $\alpha$ 恰当，势为 $K\alpha$。

这说明局部上"闭则恰当"来自空间可沿径向连续缩到一点；全局上同调正是这种收缩在整体上失败留下的信息。
\end{proof}

\begin{proposition}[正则水平集的切空间]
若 $F:M\to N$ 光滑、$y\in N$ 是正则值，则 $F^{-1}(y)$ 是余维 $\dim N$ 的嵌入子流形，且对每个 $p\in F^{-1}(y)$，$T_pF^{-1}(y)=\ker F_{*p}$。
\end{proposition}
\begin{proof}
分两步证明双向包含。先证 $T_pF^{-1}(y)\subseteq\ker F_{*p}$。任取切向量 $v\in T_pF^{-1}(y)$，由切空间定义存在光滑曲线 $\gamma:(-\delta,\delta)\to F^{-1}(y)$ 满足 $\gamma(0)=p$、$\dot\gamma(0)=v$。因 $\gamma$ 像落在水平集中，
\[
F(\gamma(t))\equiv y
\]
对一切 $t$ 成立。对 $t$ 在 $t=0$ 处求导，由链式法则
\[
F_{*p}\dot\gamma(0)
=\frac{\dd}{\dd t}\bigg|_{t=0}F(\gamma(t))
=\frac{\dd}{\dd t}\bigg|_{t=0}y=0,
\]
即 $F_{*p}v=0$，故 $v\in\ker F_{*p}$。

再证 $\ker F_{*p}\subseteq T_pF^{-1}(y)$。因 $y$ 是正则值，$F_{*p}$ 满射，由淹没正规形可取 $p$ 附近局部坐标使
\[
F(x^1,\ldots,x^m)=(x^1,\ldots,x^n).
\]
在此坐标下，水平集 $F^{-1}(y)$ 局部由
\[
x^1=y^1,\;\ldots,\;x^n=y^n
\]
定义，其余坐标 $x^{n+1},\ldots,x^m$ 自由。因此 $F^{-1}(y)$ 可参数化为
\[
(t^{n+1},\ldots,t^m)\longmapsto
(y^1,\ldots,y^n,t^{n+1},\ldots,t^m),
\]
其切空间由 $\partial_{n+1},\ldots,\partial_m$ 张成。而
\[
F_{*p}(\partial_a)=\partial_a\quad(a\le n),
\qquad
F_{*p}(\partial_a)=0\quad(a>n),
\]
故 $\ker F_{*p}=\operatorname{span}\{\partial_{n+1},\ldots,\partial_m\}=T_pF^{-1}(y)$。综合双向包含得到\eqnrefs{eq:tangent-regular-level-set}。
\end{proof}

\begin{lemma}[拉回与外微分交换]
对任意光滑映射 $F:M\to N$ 与微分形式 $\alpha\in\Omega^k(N)$，有 $\dd(F^*\alpha)=F^*(\dd\alpha)$。
\end{lemma}
\begin{proof}
目标是证明拉回与外微分可交换，即 $\dd F^*=F^*\dd$。在 $N$ 上取局部坐标 $y^a$，把任意 $k$-形式写成
\[
\alpha
=\frac1{k!}\alpha_{a_1\cdots a_k}(y)
\dd y^{a_1}\wedge\cdots\wedge\dd y^{a_k}.
\]
记 $F^a:=y^a\circ F$。由拉回定义，
\[
F^*\alpha
=\frac1{k!}(\alpha_{a_1\cdots a_k}\circ F)
\dd F^{a_1}\wedge\cdots\wedge\dd F^{a_k}.
\]
对它作用 $\dd$。由于 $\dd^2F^a=0$，外微分只作用在系数上：
\[
\begin{aligned}
\dd(F^*\alpha)
={}&\frac1{k!}
\dd(\alpha_{a_1\cdots a_k}\circ F)
\wedge\dd F^{a_1}\wedge\cdots\wedge\dd F^{a_k}\\
={}&\frac1{k!}
\bigl[(\partial_b\alpha_{a_1\cdots a_k})\circ F\bigr]
\dd F^b\wedge\dd F^{a_1}\wedge\cdots\wedge\dd F^{a_k}.
\end{aligned}
\]
另一方面，先在 $N$ 上计算
\[
\dd\alpha
=\frac1{k!}\partial_b\alpha_{a_1\cdots a_k}
\dd y^b\wedge\dd y^{a_1}\wedge\cdots\wedge\dd y^{a_k},
\]
再拉回得到完全相同的式子。因此 $\dd F^*=F^*\dd$。证明只使用链式法则与 $\dd^2=0$，所以该恒等式与任何度量或联络无关。

取 $k=1$ 且 $F(t)=(\cos t,\sin t)$ 可显式验证两端相等。
\end{proof}

\begin{theorem}[横截原像定理]
设 $S\subset N$ 是嵌入子流形，光滑映射 $F:M\to N$ 横截于 $S$。则 $F^{-1}(S)$ 为余维 $\operatorname{codim}_N S$ 的嵌入子流形，且 $T_pF^{-1}(S)=\{v\in T_pM:F_{*p}v\in T_{F(p)}S\}$。
\end{theorem}
\begin{proof}
在 $q=F(p)\in S$ 附近选适配坐标 $(y^1,\ldots,y^r,y^{r+1},\ldots,y^n)$，使
\begin{equation*}
 S\cap U=\{y^{r+1}=\cdots=y^n=0\},
 \qquad r=\dim S.
\end{equation*}
即 $S$ 的前 $r$ 个坐标自由，后 $n-r$ 个坐标为零。定义法向投影
\begin{equation*}
 \pi_\perp:N\supset U\to\RR^{n-r},
 \qquad
 \pi_\perp(y)=(y^{r+1},\ldots,y^n),
\end{equation*}
它提取后 $n-r$ 个"法向"坐标分量。

在此坐标下，
\begin{equation*}
 F^{-1}(S)=\{p\in M:F(p)\in S\}
 =\{p\in M:\pi_\perp(F(p))=0\}
 =(\pi_\perp\circ F)^{-1}(0).
\end{equation*}
因此 $F^{-1}(S)$ 是复合映射 $\pi_\perp\circ F:M\to\RR^{n-r}$ 的零水平集。

$\pi_\perp$ 的微分是到后 $n-r$ 个坐标分量的投影。复合映射的微分
\begin{equation*}
 \dd(\pi_\perp\circ F)_p=\pi_\perp\circ\dd F_p:T_pM\to\RR^{n-r}
\end{equation*}
满射，当且仅当 $\dd F_p$ 的像加上 $S$ 的切空间（即前 $r$ 个坐标方向）张成整个 $T_{F(p)}N$，这正是横截条件\eqnrefs{eq:math-dg-transversality}。

既然 $0$ 是 $\pi_\perp\circ F$ 的正则值，正则值定理给出 $F^{-1}(S)=(\pi_\perp\circ F)^{-1}(0)$ 是 $M$ 的嵌入子流形，余维数 $n-r=\operatorname{codim}_N S$，即式\eqnrefs{eq:math-dg-transverse-codimension}。切空间是 $\dd(\pi_\perp\circ F)_p$ 的核，即 $\{v\in T_pM:\dd F_p(v)\in T_{F(p)}S\}$，即\eqnrefs{eq:math-dg-transverse-tangent-space}。

横截原像定理因而不是新的局部正规形，而是把正则值定理应用于法向坐标分量。
\end{proof}

\begin{proposition}[相关向量场的流]
若 $F:M\to N$ 光滑且 $F_*X=Y\circ F$，则在两边均有定义的局部流区间内，$F\circ\Phi_t^X=\Phi_t^Y\circ F$。当 $F$ 为微分同胚时，这就是两流的共轭关系。
\end{proposition}
\begin{proof}
分两步证明 $F\circ\Phi_t^X=\Phi_t^Y\circ F$。

固定 $p\in M$，令
\[
 \eta(t)=F(\Phi_t^X(p)).
\]
这是 $F$ 沿 $X$ 的流从 $p$ 出发的轨道所得到的 $N$ 上的曲线，满足初始条件 $\eta(0)=F(p)$。

对 $\eta(t)=F(\Phi_t^X(p))$ 关于 $t$ 求导。由链式法则，
\[
 \dot\eta(t)
 =F_{*\Phi_t^X(p)}\!\left(\frac{\dd}{\dd t}\Phi_t^X(p)\right)
 =F_{*\Phi_t^X(p)}\,X_{\Phi_t^X(p)}.
\]
由相关向量场条件 $F_*X=Y\circ F$，在点 $\Phi_t^X(p)$ 处，
\[
 F_{*\Phi_t^X(p)}\,X_{\Phi_t^X(p)}
 =Y_{F(\Phi_t^X(p))}
 =Y_{\eta(t)}.
\]
因此 $\dot\eta(t)=Y_{\eta(t)}$，即 $\eta$ 是 $Y$ 的积分曲线。

另一方面，$\Phi_t^Y(F(p))$ 也是 $Y$ 的积分曲线，且满足同一初始条件：
\[
 \eta(0)=F(p)=\Phi_0^Y(F(p)).
\]
由 ODE 解的唯一性（在流的存在区间内），
\[
 \eta(t)=\Phi_t^Y(F(p)),
\]
即 $F(\Phi_t^X(p))=\Phi_t^Y(F(p))$。因 $p$ 任意，得到\eqnrefs{eq:related-flow-naturality}。
\end{proof}

\begin{proposition}[流交换子的混合二阶项]
设 $X,Y$ 为光滑向量场，流分别为 $\Phi_t,\Psi_s$。定义 $C_{t,s}=\Psi_{-s}\circ\Phi_{-t}\circ\Psi_s\circ\Phi_t$，则 $C_{t,s}^*f=f+ts[X,Y]f+O(|t|^2|s|+|t||s|^2)$。相反次序 $\Phi_t\circ\Psi_s\circ\Phi_{-t}\circ\Psi_{-s}$ 的混合项为 $-ts[X,Y]$。
\end{proposition}
\begin{proof}
将 $C_{t,s}=\Psi_{-s}\circ\Phi_{-t}\circ\Psi_s\circ\Phi_t$ 的拉回展开到 $ts$ 阶，证明最低非平凡项恰好是 Lie 括号 $[X,Y]$。

令 $\Phi_t$ 与 $\Psi_s$ 分别是 $X$ 与 $Y$ 的局部流。对任意函数 $f$，由流定义，
\begin{equation*}
 \Phi_t^*f=f+t\,Xf+O(t^2),
 \qquad
 \Psi_s^*f=f+s\,Yf+O(s^2).
\end{equation*}
等价地，拉回算子的一阶展开为
\begin{equation*}
 \Phi_{\pm t}^*=I\pm t\,X+O(t^2),
 \qquad
 \Psi_{\pm s}^*=I\pm s\,Y+O(s^2).
\end{equation*}

考虑群交换子对应的拉回组合
\begin{equation*}
 \mathcal C_{t,s}^*
 :=\Phi_t^*\Psi_s^*\Phi_{-t}^*\Psi_{-s}^*.
\end{equation*}
依次代入一阶展开，并只保留 $ts$ 阶交叉项（纯 $t$ 项与纯 $s$ 项由逆流精确抵消，余项至少含一个 $t$ 和一个 $s$）：
\begin{align*}
 \mathcal C_{t,s}^*
 &=(I+tX)(I+sY)(I-tX)(I-sY)+O(t^2s,ts^2)\\
 &=I+ts(XY-YX)+O(t^2s,ts^2)\\
 &=I+ts[X,Y]+O(t^2s,ts^2).
\end{align*}
混合项的逐项核对：前两个因子给 $tsXY$，后两个因子也给 $tsXY$；两个乘积之间的交叉项给 $-tsXY-tsYX$。合计为 $ts(XY-YX)$。纯二阶项必须连同每个流的二阶 Taylor 项一起处理，不能仅凭四个一阶多项式就声称它们消失。

这里拉回逆转映射复合的次序。若直接采用 $\Phi_t\circ\Psi_s\circ\Phi_{-t}\circ\Psi_{-s}$，其拉回依次为 $\Psi_{-s}^*\Phi_{-t}^*\Psi_s^*\Phi_t^*$，混合项因而为 $-ts[X,Y]$。

Lie 括号不是人为附加的代数运算，而是两个局部流不交换性的最低非平凡系数：若 $[X,Y]=0$，群交换子在 $ts$ 阶消失，流可交换（Frobenius 可积条件的无穷小版本）。符号由组合次序完全固定；反转群交换子次序，括号也相应变号。
\end{proof}

\begin{proposition}[Cartan 演算的闭合关系]
对光滑向量场 $X,Y$，在整个微分形式代数上有 $[\mathcal L_X,\iota_Y]=\iota_{[X,Y]}$、$[\mathcal L_X,\mathcal L_Y]=\mathcal L_{[X,Y]}$ 和 $[\mathcal L_X,\dd]=0$。插入算子满足 $[\iota_X,\iota_Y]_{\rm gr}=0$。
\end{proposition}
\begin{proof}
先在函数与一形式上验证
\[
[\mathcal L_X,\iota_Y]=\iota_{[X,Y]}.
\]
对一形式 $\alpha$，
\[
\begin{aligned}
([\mathcal L_X,\iota_Y]\alpha)
&=X(\alpha(Y))-(\mathcal L_X\alpha)(Y)\\
&=X(\alpha(Y))-
\bigl[X(\alpha(Y))-\alpha([X,Y])\bigr]\\
&=\alpha([X,Y]).
\end{aligned}
\]
两边都是次数 $-1$ 的导子，所以在整个外代数上一致，得到式\eqnrefs{eq:cartan-lie-iota}。再用分次 Jacobi 恒等式与
\[
\mathcal L_X=[\dd,\iota_X]_{\rm gr}
\]
可得
\[
\begin{aligned}
[\mathcal L_X,\mathcal L_Y]
&=[\mathcal L_X,[\dd,\iota_Y]_{\rm gr}]\\
&=[\dd,[\mathcal L_X,\iota_Y]]_{\rm gr}\\
&=[\dd,\iota_{[X,Y]}]_{\rm gr}\\
&=\mathcal L_{[X,Y]},
\end{aligned}
\]
即\eqnrefs{eq:cartan-lie-lie}。所以 Lie 括号、流与微分形式代数三者并不是并列工具，而是由同一分次交换子结构闭合起来。

在 $\RR^3$ 上显式验证 Cartan 恒等式。取 $X=\partial_x$，$\alpha=\dd y$。则 $\iota_X\alpha=0$，$\mathcal L_X\alpha=\dd(\iota_X\alpha)+\iota_X\dd\alpha=0$（因 $\dd\alpha=0$）。另一方向，$\mathcal L_X\dd y=\dd(\mathcal L_X y)=\dd(0)=0$，验证 $[\mathcal L_X,\dd]=0$。取 $Y=x\partial_y$，$[X,Y]=\partial_y$，则 $[\mathcal L_X,\iota_Y]\dd y=\iota_{[X,Y]}\dd y=1$，直接计算 $\mathcal L_X(\dd y(Y))-\iota_Y(\mathcal L_X\dd y)=\mathcal L_X(x)-0=1$，验证\eqnrefs{eq:cartan-lie-iota}。

这一分次微分结构在多处反复出现：Bianchi 恒等式 $\dd F=\dd^2A=0$ 给出齐次 Maxwell 方程，不可压 Euler 的涡度方程 $\partial_t\omega+\mathcal L_v\omega=0$ 是 Cartan 恒等式的特例，Maurer--Cartan 方程则是 $\dd^2=0$ 的非线性推广。

故外微分与 Hodge 对偶把电磁场强、Bianchi 恒等式与 Maxwell 方程 $\dd *F=J$ 写成统一语言，其核心 $\dd^2=0$ 与 Cartan 恒等式根源即外代数上的分次微分结构。
\end{proof}

\begin{theorem}[Cartan 恒等式]
对任意光滑向量场 $X$ 与任意次数的微分形式，Lie 导数满足 $\mathcal L_X=\iota_X\dd+\dd\iota_X$。该恒等式不需要度量或联络。
\end{theorem}
\begin{proof}
目标是证明 Cartan 恒等式 $\mathcal L_X=\iota_X\dd+\dd\iota_X$。定义次数为零的算子
\begin{equation*}
\mathscr D_X:=\mathcal L_X-\iota_{X}\dd-\dd\iota_{X}.
\end{equation*}
对任意函数 $f$，由 $\iota_{X}f=0$ 与\eqnrefs{eq:math-dg-differential-function}可得
\begin{align*}
\mathscr D_Xf
&=\mathcal L_X f-\iota_{X}\dd f-\dd(\iota_{X}f)\\
&=X(f)-\iota_{X}\dd f-\dd 0\\
&=X(f)-\dd f(X)=0.
\end{align*}
Lie 导数是次数零导子，而 $\iota_{X}\dd+\dd\iota_{X}$ 也是次数零导子：对任意形式 $\omega,\eta$，
\begin{align*}
(\iota_{X}\dd+\dd\iota_{X})(\omega\wedge\eta)
&=\iota_{X}\dd(\omega\wedge\eta)+\dd\iota_{X}(\omega\wedge\eta)\\
&=\iota_{X}(\dd\omega\wedge\eta+(-1)^{|\omega|}\omega\wedge\dd\eta)\\
&\quad+\dd(\iota_{X}\omega\wedge\eta+(-1)^{|\omega|}\omega\wedge\iota_{X}\eta),
\end{align*}
按 $\iota_{X}$ 与 $\dd$ 各自的反导子 Leibniz 规则展开后，非齐次项相消，剩余项恰为
\begin{equation*}
(\iota_{X}\dd+\dd\iota_{X})\omega\wedge\eta
+\omega\wedge(\iota_{X}\dd+\dd\iota_{X})\eta.
\end{equation*}
因此 $\mathscr D_X$ 满足普通 Leibniz 规则。又由\eqnrefs{eq:math-dg-synthesis-lie-d-commute}和 $\dd^2=0$，
\begin{align*}
[\mathscr D_X,\dd]
&=[\mathcal L_X,\dd]-[\iota_{X}\dd,\dd]-[\dd\iota_{X},\dd]\\
&=0-\iota_{X}\dd^2+\dd\iota_{X}\dd-\dd\iota_{X}\dd+\dd^2\iota_{X}=0.
\end{align*}
于是对任意函数 $f$，
\begin{equation*}
\mathscr D_X(\dd f)=\dd(\mathscr D_X f)=\dd 0=0.
\end{equation*}
局部上任意微分形式都由光滑函数与坐标一形式 $\dd x^i$ 通过外积生成，而 $\mathscr D_X$ 既是导子又与 $\dd$ 交换，故对单项式 $f_0\,\dd f_1\wedge\cdots\wedge\dd f_k$ 有
\begin{align*}
\mathscr D_X(f_0\,\dd f_1\wedge\cdots\wedge\dd f_k)
&=(\mathscr D_X f_0)\,\dd f_1\wedge\cdots\wedge\dd f_k\\
&\quad+\sum_{j=1}^k f_0\,\dd f_1\wedge\cdots\wedge\mathscr D_X(\dd f_j)\wedge\cdots\wedge\dd f_k\\
&=0.
\end{align*}
所以导子 $\mathscr D_X$ 在整个 $\Omega^\bullet(M)$ 上恒为零。故
\begin{equation}\label{eq:math-dg-synthesis-cartan}
\boxed{\mathcal L_X=\iota_{X}\dd+\dd\iota_{X}.}
\end{equation}
\end{proof}
